\documentclass{./texmf/krvntu}

\usepackage{cmap}
\usepackage[T2A]{fontenc}
\usepackage[utf8]{inputenc}
\usepackage[english,ukrainian]{babel}
\usepackage{./texmf/pscyr}

\input{./texmf/t2aftm.fd}
\input{./texmf/t2afer.fd}

\pdfmapfile{+pscyr.map}

\renewcommand{\rmdefault}{ftm}
\renewcommand{\sfdefault}{fer}

% Визначення команд форматування
\newcommand{\term}[1]{\textbf{#1}}
\newcommand{\important}[1]{\textbf{\textit{#1}}}
\newcommand{\example}[1]{\textit{#1}}

\begin{document}

\section{Тест шрифтів PSCyr}

Це звичайний текст українською мовою.

\textbf{Це жирний текст українською мовою.}

\textit{Це курсивний текст українською мовою.}

\textbf{\textit{Це жирний курсивний текст українською мовою.}}

\section*{Перевірка команд}

\term{Ключовий термін} — використовує жирний шрифт.

\important{Важливе поняття} — використовує жирний курсив.

\example{Приклад або іноземний термін} — використовує курсив.

\end{document}
